\section{Die zentrale Formel}

\textbf{[K]} Die Gravitationskonstante ist in dieser Theorie keine unabhängige
Größe, sondern folgt aus $\xi$ und einer Masse:
\[
  G = \frac{\xi^2}{4\,m_\text{char}} \qquad\text{(natürliche Einheiten)}.
\]

Das ist eine \emph{Herleitung}, keine Brücke. G wird nicht als Einheitenfaktor
absorbiert (das wäre A080), sondern als Funktion des einen Parameters und einer
charakteristischen Masse ausgedrückt. Der Unterschied zu A130 ist
entscheidend: dort war $\alpha$ am Ende kalibriert; hier ist die
\emph{Struktur} von G aus $\xi$ bestimmt.

\textbf{[B]} Die äquivalente Form
\[
  \xi = 2\sqrt{G\,m_\text{char}}
\]
zeigt die Aussage von der anderen Seite: $\xi$, die tetraedrische
Packungsgröße aus A020, \emph{ist} bis auf den Faktor die Wurzel aus dem Produkt
von Gravitationsstärke und charakteristischer Masse. Gravitation und Geometrie
sind damit nicht zwei Dinge, sondern zwei Lesarten von $\xi$.

Dies ist nur in natürlichen Einheiten (A085) eine Aussage über Struktur. In SI
wäre es eine Gleichung zwischen inkommensurablen Einheiten.

\textbf{[K]} Zur Einordnung: in natürlichen Einheiten ist $G = 1$, genau wie
$c = \hbar = k_B = \alpha = 1$ (A085, A130). Der SI-Wert $6{,}674\cdot10^{-11}$
ist Buchhaltung. Die Formel $G = \xi^2/(4m)$ sagt daher nicht, welchen nativen
Wert $G$ hat -- der ist Eins --, sondern wie sich die \emph{SI-Kleidung} von $G$
aus $\xi$ und einer Masse zusammensetzt. Das ist dieselbe Buchhaltung wie bei
$\alpha$, nur für eine dimensionsbehaftete statt einer dimensionslosen Größe.

\section{Warum das eine besonders direkte Herleitung ist}

\textbf{[K]} Wie bei $\alpha$ (A130) steht auch hier keine Zirkularität im Kern
der Formel -- bei $\alpha$ ist die anfängliche Zirkularität durch die
Zwei-Wege-Überbestimmung von $E_0$ gebrochen, bei G ist sie von vornherein nicht
vorhanden: $G = \xi^2/(4m)$ verwendet $G$ an keiner Stelle, um $G$ zu bestimmen.
Der einzige Unterschied ist die Unmittelbarkeit -- $\alpha$ braucht dafür das
Zusammentreffen zweier Wege, G braucht es nicht, weil die Formel $G$ schlicht
nicht enthält. Eine versteckte Kalibrierung sitzt auch hier nicht im Kern,
sondern erst in der SI-Kleidung ($E_\text{char}$, siehe unten). Die Formel sagt:
wähle die charakteristische Masse -- etwa $m_e$ --, und die Gravitationsstärke
folgt aus $\xi$.

\textbf{[B]} Das ist die eigentliche Auflösung des Gravitationsproblems in
dieser Theorie, und sie ist stärker als die in A140 genannte. A140 sagt: das
QG-Problem stellt sich nicht, weil die Geometrie vorgeordnet ist. Das ist ein
\emph{negatives} Argument. A145 sagt: G hat einen \emph{positiven} Ausdruck aus
$\xi$. Das negative Argument erklärt, warum man nicht quantisieren muss; das
positive erklärt, woher der Zahlenwert kommt. Das macht G zur direktesten
$\xi$-Manifestation des Korpus -- nicht zur genauest geprüften (das ist die
Leptonleiter A110 bzw.\ die $\alpha$-Überbestimmung A130), aber zur
unmittelbarsten in der Form.

\section{Die Dimensionsbrücke E\_char}

\textbf{[K]} Die Grundformel $G=\xi^2/(4m)$ ist dimensional noch nicht
vollständig: in natürlichen Einheiten hat $G$ die Dimension $[E^{-2}]$, die
rechte Seite aber $[E^{-1}]$. Es fehlt ein Faktor $1/E_\text{char}$:
\[
  G_\text{nat} = \frac{\xi_0^2}{4\,m_e}\cdot\frac{1}{E_\text{char}},
  \qquad E_\text{char}\approx 28{,}4 .
\]

\textbf{[K]} Der Korpus gibt für $E_\text{char}$ eine strukturelle Zerlegung:
\[
  E_\text{char} = E_0\cdot\Bigl(\tfrac43\Bigr)^2\cdot\frac{\pi}{\sqrt2}\cdot
  K_\text{frak} = 7{,}40\cdot 1{,}778\cdot 2{,}221\cdot 0{,}986 = 28{,}8 ,
\]
gegen den Zielwert $28{,}4$, also $+1{,}5\,\%$.

\textbf{Die Einordnung -- und sie entlastet.} $E_\text{char}$ trägt zwei Größen,
die anderswo nicht rein geometrisch sind: $E_0$ (A130) und $K_\text{frak}$
(A040). Der Faktor $(4/3)^2$ ist geometrisch, $\pi/\sqrt2$ ebenfalls; $E_0$ und
$K_\text{frak}$ sind es nicht. Aber das betrifft allein den \emph{SI-Zahlenwert}
von $G$ -- und der ist Buchhaltung. Wie genau $E_\text{char}$ die $28{,}4$
trifft, ist eine Aussage über die Skalenumrechnung, nicht über die Struktur.
Solange man mit Verhältnissen rechnet, ist der Messwert von $G$ unwichtig; die
Formelzusammenhänge sind in SI und im Standardmodell bekannt und bestätigt. Die
\emph{Grundstruktur} $G = \xi^2/(4m)$ ist davon unberührt und ist die eigentliche
Aussage.

Sauber getrennt: die Aussage \enquote{G ist proportional zu $\xi^2$ und
umgekehrt proportional zu einer Masse} ist \textbf{[K]}. Der genaue
Zahlenwert über $E_\text{char}$ ist \textbf{[B]} mit kalibriertem Anteil.

\section{Dieselbe Diagnose wie bei $\alpha$, komplexere Kleidung}

\textbf{[K]} G besteht denselben Test, an dem sich in A135 entscheidet, ob eine
Größe grundlegend ist: \emph{lässt sie sich durch widerspruchsfreie
Einheitenwahl auf Eins setzen?} Für G ist die Antwort ja -- in Planck-Einheiten
gilt $G = c = \hbar = 1$. Nach der Diagnose aus A135 heißt das: G ist kein
Inhalt der Theorie, sondern ein Umrechnungsfaktor. Es steht damit auf derselben
Seite wie $c$, $\hbar$, $k_B$ (physikalisch auf Eins über die Dualität) und wie
$\alpha$ (per Normierung auf Eins) -- alles Konvention. Der SI-Wert
$6{,}674\cdot10^{-11}$ ist das, was diese Konvention verbirgt, nicht ein
grundlegender Zahlenwert.

\textbf{[K]} Und wie bei $\alpha$ liegt der Inhalt nicht in diesem Zahlenwert,
sondern in $\xi$. Die blanke Münze ist $\xi$; $G$ ist -- in SI-Licht gesehen --
dieselbe Münze, nur eingekleidet. Das ist der Sinn von $G = \xi^2/(4m_\text{char})$:
nicht ein zweiter Freiheitsgrad, sondern $\xi$ in dimensionsbehafteter Gestalt.
Äquivalenter Informationsgehalt, nicht gleicher Rang -- genau die Auflösung der
scheinbaren Spannung aus A135.

\textbf{[K]} Der Unterschied zu $\alpha$ ist allein die \emph{Reichhaltigkeit
der Kleidung}, und er hat einen dimensionalen Grund. In natürlichen Einheiten
ist $\alpha$ dimensionslos, $[E^0]$; ihre SI-Gestalt braucht daher genau eine
Skala:
\[
  \alpha_\text{SI} = \xi\cdot E_0^2 \qquad\text{(eine Skala, quadriert).}
\]
$G$ dagegen hat in natürlichen Einheiten die Dimension $[E^{-2}]$
(es ist $M_\text{Planck}^{-2}$) und bindet damit Masse an Länge an Zeit. Seine
SI-Gestalt kann deshalb nicht mit einer Skala auskommen; sie braucht eine Masse
\emph{und} die Zwischenskala $E_\text{char}$:
\[
  G_\text{nat} = \frac{\xi^2}{4\,m_\text{char}}\cdot\frac{1}{E_\text{char}}.
\]
$E_\text{char}$ ist selbst mehrschichtig -- $E_0\cdot(4/3)^2\cdot(\pi/\sqrt2)\cdot
K_\text{frak}$ (siehe oben) --, während $\alpha$ mit dem einen $E_0$ auskommt.
Das ist die \enquote{komplexere SI-Verkleidung}: nicht mehr Inhalt, sondern mehr
Schichten desselben Inhalts, erzwungen durch die reichere Dimension von G.

\section{Was hier zählt: die exakte Rückrechenbarkeit}

\textbf{[K]} Die tragende Aussage dieses Dokuments ist nicht, wie genau die
$0{,}5\,\%$ getroffen werden, sondern dass G in natürlichen Einheiten auf Eins
setzbar ist -- und daraus folgt etwas Exaktes. Ist G eine Konvention (Eins in
Planck-Einheiten), dann ist sein SI-Wert kein unabhängiger Zahlenwert, sondern
\emph{vollständig} durch die Umrechnung festgelegt:
\[
  G_\text{SI} = \frac{\hbar c}{M_\text{Planck}^2} .
\]
Gibt man $\xi$ und die eine absolute Skala, so ist $G_\text{SI}$ damit
\emph{exakt} bestimmt -- ohne freien Parameter, ohne Näherung. Das ist dieselbe
Struktur wie $\alpha_\text{SI} = \xi\,E_0^2$: die dimensionsbehaftete Konstante
ist die SI-Aussprache von $\xi$ plus Skala, und die Rückrechnung ist exakt, nicht
approximativ.

\textbf{[K]} Davon strikt zu trennen ist die $E_\text{char}$-Zerlegung. Das exakte
$E_\text{char}$ -- dasjenige, das $G_\text{SI}$ trifft -- ist mit $\xi$ und der
Skala bereits festgelegt. Was $0{,}5\,\%$ danebenliegt, ist nicht diese exakte
Größe, sondern der \emph{Versuch, sie durch eine geschlossene Formel} aus
bekannten Größen auszudrücken ($E_0\cdot(4/3)^2\cdot(\pi/\sqrt2)\cdot
K_\text{frak}$). Die Abweichung misst die Güte dieser \emph{analytischen
Näherung der Kleidung}, nicht die Rückrechenbarkeit selbst. Die Rückrechnung ist
exakt; nur ihre Abkürzung in einen einfachen Ausdruck ist es noch nicht.

\textbf{[SETZUNG]} Damit steht die Priorität richtig: Was G von einer echten
fundamentalen Konstante unterscheidet, ist die Setzbarkeit auf Eins und die
daraus folgende exakte Rückrechenbarkeit -- und die ist gezeigt. Der Rest der
$E_\text{char}$-Näherung ist eine Frage der geschlossenen Form, nicht der
Struktur, und für die tragende Aussage ohne Gewicht.

\section{Die Hierarchie der Skalen}

\textbf{[B]} Die Formel ordnet drei Energieskalen:
\[
  E_0 = 7{,}40\ \text{MeV}\ \ll\ E_\text{char} = 28{,}4\ \ll  E_{T0} = \frac{1}{\xi_0} = 7500 .
\]

$E_0$ ist die elektromagnetische Skala (A130), $E_{T0}=1/\xi$ die
fundamentale Skala der Theorie, $E_\text{char}$ die Zwischenskala, auf der die
Gravitation ankoppelt. Dass die Gravitation bei einer \emph{Zwischenskala}
sitzt, ist die strukturelle Aussage hinter ihrer Schwäche: sie koppelt weit
unterhalb der fundamentalen Skala.

\textbf{[S]} Ob diese Hierarchie die beobachtete Schwäche der Gravitation
\emph{quantitativ} erklärt -- die vierzig Größenordnungen gegenüber dem
Elektromagnetismus --, ist damit nicht gezeigt. Es ist eine Richtung, keine
Rechnung.

\section{Prüfskript}

\begin{itemize}
  \item $G=\xi^2/(4m_e)$ und die Rückrichtung $\xi=2\sqrt{Gm}$; die
    Dimensionslücke und ihre Schließung durch $E_\text{char}$; die
    strukturelle Zerlegung von $E_\text{char}$ mit Ausweis des kalibrierten
    Anteils; und die Skalenhierarchie:
    \texttt{python/A\_Serie\_Skripte/a145\_gravitationskonstante.py}
\end{itemize}


\section{Zeichenerklärung}

Symbole die in der gesamten A-Serie verwendet werden, sind in A010 erklärt.
Spezifisch für dieses Dokument:

\begin{center}
\renewcommand{\arraystretch}{1.3}
{\small\begin{tabular}{lp{9cm}}
\toprule
Symbol & Bedeutung \\
\midrule
$m_\text{char}$ & Charakteristische Masse; Grenzfall der Selbstkonsistenz-
  bedingung, $G = \xi^2/(4m_\text{char})$ \\
$G$ & Newtonsche Gravitationskonstante \\
\bottomrule
\end{tabular}}
\renewcommand{\arraystretch}{1.0}
\end{center}

\section{Was offen bleibt}

\textbf{Erstens ist die charakteristische Masse nicht eindeutig festgelegt.} Der
Korpus verwendet $m_e$; warum das Elektron und nicht eine andere Masse die
charakteristische ist, ist nicht begründet.

\textbf{Zweitens trägt $E_\text{char}$ einen kalibrierten Anteil} (siehe oben).
Die Grundstruktur $\xi^2/(4m)$ ist davon unberührt, der genaue Zahlenwert
nicht. \emph{Dieser Rest ist ausdrücklich kein Messeffekt.} Es liegt nahe, die
$\sim0{,}5\,\%$ auf die berüchtigte Messschwäche von G schieben zu wollen -- doch
das Gegenteil ist der Fall: G ist auf $2\cdot10^{-5}$ (rund $0{,}002\,\%$)
gemessen, und selbst die große Streuung der Einzelexperimente untereinander
erreicht nur $\sim0{,}05\,\%$. Der Theorie-Rest ist damit rund $200$-mal größer
als die Messunsicherheit und etwa $10$-mal größer als die experimentelle
Streuung; die Eingänge der Verkleidung ($E_0$ über die Massen auf $\sim10^{-8}$,
$K_\text{frak}$ abgeleitet, die Geometriefaktoren exakt) tragen ihn ebenso wenig.
Der Rest ist also strukturell -- ein Rest der $E_\text{char}$-Näherung -- und
muss als solcher benannt bleiben; ihn der Messungenauigkeit zuzuschreiben wäre
dieselbe Ziel-zuerst-Zuordnung, vor der A105 und A130 warnen.

\textbf{Drittens ist der Rest mit einer Standardmodell-Annahme verschränkt --
der Massendefinition.} In die Verkleidung gehen \emph{gemessene} Leptonmassen
ein ($E_0=\sqrt{m_e m_\mu}$, $m_\text{char}=m_e$), und \enquote{Masse} ist im
Standardmodell schemaabhängig: die Messwerte sind \emph{Polmassen}, während eine
rein geometrische Konstruktion ebenso gut einer laufenden (MS-quer-)Masse
entsprechen könnte. Der Unterschied beider Schemata ist von der Ordnung
$\alpha/\pi \approx 0{,}23\,\%$ pro Masse; da G über
$G\sim 1/(m_e^{3/2} m_\mu^{1/2})$ mit der Gesamt-Massenpotenz $2$ von den Massen
abhängt, pflanzt sich das zu $\sim 0{,}46\,\%$ auf G fort -- \emph{derselben
Größenordnung wie der Rest}. Das entlastet den geometrischen Kern in beide
Richtungen und belastet ihn in keine: es zeigt, dass die $0{,}5\,\%$
\emph{nicht} sauber als geometrischer Defekt gelesen werden dürfen, weil ein Teil
davon eine Konvention auf der SM-Eingangsseite sein kann; es zeigt aber
\emph{nicht}, dass der Schema-Effekt die Lücke schließt. Das wäre erst gezeigt,
wenn man die Pol-zu-laufend-Verschiebung für genau diese Massen und diese Formel
mit Vorzeichen und Koeffizient \emph{vorwärts} rechnet -- ein Ordnungs-Argument
ist keine Herleitung. Festzuhalten bleibt allein: der Rest sitzt zwischen
geometrischer Näherung und Massenschema, und beides ist von gleicher Ordnung.

\textbf{Viertens fehlt weiterhin die Gravitationsdynamik.} A145 gibt den
\emph{Wert} von G aus $\xi$; es gibt nicht die Feldgleichungen, aus denen sich
das Abstandsgesetz und die Bewegung ergäben. Das bleibt der große offene Punkt,
den schon A140 nennt.

\section{Ersetzt}

Dieses Dokument nimmt auf: Dok.~012 (Gravitationskonstante, Herleitungsteil),
Dok.~127 (Gravitation, Planck-Perspektive). Es ergänzt A140 um die positive
Herleitung; A140 behält die Brücken- und QG-Argumentation.
Eingearbeitet: P40. Neu: die Messung von G am Eins-Test (Planck-Einheiten $G=c=\hbar=1$) und der Kleidungsvergleich zu $\alpha$ (dimensionaler Grund der reicheren SI-Verkleidung).

\section{Verwendung von KI-Werkzeugen}

Dieses Dokument ist unter Verwendung KI-gestützter Werkzeuge entstanden. Die
Fragestellungen, die Setzungen und alle inhaltlichen Entscheidungen stammen vom
Autor; die Werkzeuge formulieren aus, rechnen nach und erstellen Prüfskripte.
Die Verantwortung für Inhalt und Fehler liegt vollständig beim Autor. Der
ausführliche Wortlaut steht in A010.
